\begin{ExerciseBlock}
\ExerciseHeading{07}{Anharmonische Schwingung}{free_oscillations_anharmonic_trap}
In einer eindimensionalen Ionenfalle wirkt auf ein Ion der Masse $m$ die Kraft
\[
 F(x)=-F_0 x\,\mathrm e^{-(x/\ell)^2}.
\]
Dabei ist $\ell>0$ eine charakteristische Länge und $F_0>0$ hat die Einheit einer Federkonstante.
\begin{enumerate}
\item Skizzieren Sie $F(x)$ und identifizieren Sie die Ruhelage.
\item Stellen Sie die Bewegungsgleichung auf. Handelt es sich um eine harmonische Schwingung?
\item Nähern Sie $F(x)$ für kleine Auslenkungen durch die Tangente $F_\mathrm T(x)$ in der Ruhelage an.
\item Stellen Sie die Bewegungsgleichung für $F_\mathrm T(x)$ auf. Begründen Sie, ob die genäherte Bewegung harmonisch ist, und geben Sie gegebenenfalls ihre Kreisfrequenz an.
\end{enumerate}
\end{ExerciseBlock}
\UseExerciseSolution{free_oscillations_anharmonic_trap}
